\clearpage
\section{Privileged Execution Model}\label{page:privileged-execution-model}

All exceptions, interrupts, NMI, and explicit synchronous traps enter supervisor execution through the architectural event mechanism. \texttt{SYSCALL} is the explicit \texttt{SYSTEM\_CALL} event (\texttt{EXC:0x20}); it uses the common EPC/ECS/EDS entry and returns with \texttt{ERET}.

\subsection{Privilege and Event State}
STATUS.PM is zero in user mode and one in supervisor mode. STATUS.EA indicates active event processing. STATUS.EDEPTH in bits 9..6 is the active event depth, and STATUS.UO in bit 10 records that the outer event originated in user mode. PM, EA, EDEPTH, and UO are hardware managed; WRSTATUS must preserve them. User execution always observes PM=EA=EDEPTH=UO=0.

IE enables maskable interrupts, TF requests single-step tracing, RF suppresses one trace boundary, and NI inhibits nested NMI admission. RDCR, WRCR, translation-state changes, event-control changes, and page-table context switches require supervisor privilege unless an instruction explicitly states otherwise.

\subsection{User Context Bank}
The user context bank consists of UPC, USP, UCS, UDS, USS, UCTL, and UINFO. UCTL contains saved FLAGS in bits 15..0, saved STATUS in bits 31..16, and V in bit 32. UINFO contains the 32-bit event code; its upper half is zero. A saved user STATUS image has PM, EA, EDEPTH, and UO clear.

Software stages a context by writing UPC, USP, UCS, UDS, USS, and UINFO before writing UCTL.V=1. That final write validates the complete bank. The bank belongs to the interrupted user thread and must be saved with its supervisor context across preemption or migration.

\subsection{Initial Event Stacks and Payloads}
A first user-origin \texttt{SYSTEM\_CALL} selects SSS:SSP, a maskable interrupt selects ISS:ISP, and an exception or NMI selects FSS:FSP. BASIC events, including \texttt{SYSTEM\_CALL}, have no payload and perform no stack-memory access. ERROR events store 16 bytes; PAGE\_FAULT and AUXILIARY events store 32 bytes. The event code in UINFO defines the payload shape.

A supervisor-origin event and every nested event push a complete event frame on the current SS:SP. They do not promote to another configured event stack. The saved SS:SP in the frame identifies the interrupted stack.

\subsection{Event Entry and Failure}
User-origin entry atomically saves the user context bank, optionally stores the payload, loads EPC/ECS/EDS and the selected initial stack, sets PM and EA, sets EDEPTH to one and UO to one, and clears IE, TF, and RF. Supervisor and nested entry save the prior context in the full frame and increment EDEPTH while preserving UO.

Delivery failure selects DSS:DSP and attempts \texttt{DOUBLE\_FAULT}. The hidden current-DFA bit is set when that entry commits. Failure while current DFA is set, or failure of the double-fault attempt, enters shutdown. A full frame saves the prior DFA state in FRAME\_CONTROL.SAVED\_DFA.

\subsection{ERET Routing and Atomicity}
ERET classifies its source before any stack access. PM=1, EA=0, EDEPTH=0, UO=0, and UCTL.V=1 performs direct initial user entry from the bank. PM=1, EA=1, EDEPTH=1, UO=1, and UCTL.V=1 returns the outer user event from the bank. Other supervisor states with EA=1 and nonzero EDEPTH restore a full frame from the current SS:SP. Every other combination raises \texttt{INVALID\_CONTROL\_STATE}.

A U-bank return validates UCTL, UINFO, the user segment images, USP, and the executable UPC before committing. It restores user execution, clears UCTL.V, EDEPTH, UO, EA, and current DFA, and never reads the stack. A stacked return validates the entire frame, requires saved STATUS.EDEPTH+1 to equal the current depth and saved UO to match the current chain, then restores STATUS and SAVED\_DFA atomically. A failed validation or target probe leaves the current context and return source unchanged.

Normative instruction-local behavior is defined by \hyperref[instr:syscall]{\texttt{SYSCALL}} and \hyperref[instr:eret]{\texttt{ERET}}.\par
